\documentclass[11pt]{article}
\usepackage[margin=1in]{geometry}
\usepackage{amsmath,amssymb}
\usepackage{booktabs}
\usepackage[dvipsnames]{xcolor}
\usepackage{hyperref}
\hypersetup{colorlinks=true,linkcolor=blue!50!black,urlcolor=blue!50!black,citecolor=blue!50!black}

\title{A Fixed P3b Endpoint Carries Over to an\\
Independent Auditory Oddball Dataset:\\
A Preregistered Cross-Paradigm Robustness Test}
\author{David Lindgreen}
\date{October 9, 2026 \\[2pt] \small Version 0.1.0 \quad Preprint, not peer reviewed}

\emergencystretch=3em
\begin{document}
\maketitle

\begin{abstract}
The P3b is among the most studied event-related potentials, but a measurement chosen after seeing data can find an effect
almost anywhere on the scalp. We therefore fixed one endpoint in advance, the mean voltage at Pz from 300 to 600\,ms after
stimulus onset recommended by ERP CORE~\cite{kappenman2021}, with its contrast direction, participant-level unit of
inference, artifact threshold and stopping rule, before inspecting any amplitude in an independent dataset: the auditory
three-stimulus oddball task of OpenNeuro \texttt{ds003061}~\cite{delorme2021}. All 13 participants showed a positive
correct-target-minus-correct-standard contrast; the mean was $+5.65\,\mu$V (95\% interval 4.83 to 6.48),
$t(12)=15.00$, two-sided $p=3.9\times10^{-9}$, Cohen's $d_z=4.16$, unchanged across three artifact thresholds. Recomputing every
statistic from the 13 stored participant contrasts reproduces the report exactly and adds small-sample quantities: an exact
sign-flip $p=2.4\times10^{-4}$, a noncentral-$t$ interval for $d_z$ of 2.42 to 5.88, and a 95\% prediction interval for a new
participant of $+2.58$ to $+8.73\,\mu$V. This is cross-paradigm robustness evidence, not a direct replication.
\end{abstract}

\section{Question}
Does the classic P3b target enhancement survive when the measurement is fixed in advance and applied to an independent public
EEG dataset with a different sensory modality? ERP CORE used an active visual oddball task; the external dataset used an
active auditory oddball task~\cite{pernet2019}. The test is therefore a robustness check of a fixed endpoint, not a literal replication.

\section{Preregistered design}
The endpoint was frozen before any EEG amplitude was inspected: electrode Pz, mean voltage 300 to 600\,ms post-stimulus, contrast
correct targets minus correct standards, one contrast per participant (trials are pooled within participant, never treated as
independent people), a $150\,\mu$V peak-to-peak artifact rule, and a minimum of 20 correct targets and 100 correct standards per
run. The dataset metadata was audited at a pinned commit without inspecting signals: all 39 channel tables contain Pz, sampling is
256\,Hz throughout, and there are three runs per participant. One visibly truncated recording (\texttt{sub-012} run 1) failed the
frozen minimum before any signal analysis and was excluded; the participant remained through runs 2 and 3. The 38 eligible runs
contributed 3{,}606 accepted target and 19{,}481 accepted standard epochs.

\section{Confirmatory result}
\begin{table}[h]
\centering\small
\begin{tabular}{lr}
\toprule
Quantity & Estimate\\
\midrule
Participants & 13\\
Mean target-minus-standard contrast & $+5.654\,\mu$V\\
Median contrast & $+5.570\,\mu$V\\
95\% interval for the mean & $+4.833$ to $+6.476\,\mu$V\\
Two-sided one-sample test & $t(12)=15.004$, $p=3.88\times10^{-9}$\\
Cohen's $d_z$ & 4.161\\
Positive participant contrasts & 13/13 (range $+3.22$ to $+8.72\,\mu$V)\\
Wilcoxon sensitivity & $W=0$, $p=0.000244$\\
\bottomrule
\end{tabular}
\caption{Preregistered confirmatory result at Pz, 300--600\,ms.}
\end{table}
The prespecified artifact-threshold sensitivity leaves the conclusion unchanged: 100\,$\mu$V gives $+5.827$ (4.893 to 6.761);
150\,$\mu$V (confirmatory) $+5.654$; 200\,$\mu$V $+5.738$ (4.825 to 6.652), with all 13 participants retained in each.

\section{Post-hoc recomputation from participant contrasts}
Because the endpoint is one number per participant, every reported statistic can be recomputed exactly from the 13 stored contrasts,
without the raw EEG. We did so as an independent check, and added quantities suited to $n=13$. Not preregistered; all reported.
\begin{table}[h]
\centering\small
\begin{tabular}{lr}
\toprule
Quantity & Value\\
\midrule
Largest difference from the reported mean, $t$, $d_z$ and interval & 0\\
Hedges' small-sample $g_z$~\cite{hedges1981} & 3.896\\
Noncentral-$t$ 95\% interval for $d_z$ & 2.424 to 5.883\\
Exact sign-flip permutation $p$ (two-sided, $2^{13}$ patterns) & 0.000244\\
Leave-one-participant-out mean range & $+5.399$ to $+5.858\,\mu$V\\
Leave-one-out smallest $t$ / largest $p$ & 13.77 / $2.8\times10^{-8}$\\
95\% prediction interval, a new participant & $+2.58$ to $+8.73\,\mu$V\\
\bottomrule
\end{tabular}
\caption{Recomputation from the participant-level contrasts (\texttt{results/summary.json}).}
\end{table}

\section{Discussion and limits}
The effect is large, present in every participant, and not driven by any single person or artifact threshold. The prediction interval
excludes zero under a normality assumption, a statement about a new draw from this population rather than about all auditory
oddball studies. What the result supports is narrow: a fixed posterior 300--600\,ms target enhancement appears in an independent,
small auditory dataset despite a change of modality and paradigm, which a post-hoc search of the scalp could not demonstrate. It does
not show that visual and auditory oddball tasks engage identical processes, localize a generator, or support any diagnostic or
individual-prediction use. The sample is 13 people. The recomputation here starts from the stored participant contrasts; it did not
re-run signal processing from the raw EEG, whose download and preprocessing are specified in the repository and verified there
by checksums.

\section*{Reproducibility and AI assistance}
Code, protocol, metadata audit, results and tests:\\
\url{https://github.com/lindgreendavid/neuro-signal-lab}\\
Interactive laboratory:\\
\url{https://lindgreendavid.github.io/neuro-signal-lab/}\\
The analysis code, tests and drafts of this text were produced with substantial assistance from an AI system (Claude, Anthropic);
the author is responsible for the content.

\begin{thebibliography}{9}
\bibitem{kappenman2021} E.~S. Kappenman, J.~L. Farrens, W.~Zhang, A.~X. Stewart, S.~J. Luck, ERP CORE: An open resource for
human event-related potential research, \emph{NeuroImage} \textbf{225}, 117465 (2021).
\bibitem{delorme2021} A.~Delorme, EEG data from an auditory oddball task, OpenNeuro \texttt{ds003061} v1.1.0,
doi:10.18112/openneuro.ds003061.v1.1.0.
\bibitem{pernet2019} C.~R. Pernet et al., EEG-BIDS, an extension to the brain imaging data structure for electroencephalography,
\emph{Sci. Data} \textbf{6}, 103 (2019).
\bibitem{hedges1981} L.~V. Hedges, Distribution theory for Glass's estimator of effect size and related estimators,
\emph{J. Educ. Stat.} \textbf{6}, 107--128 (1981).
\end{thebibliography}
\end{document}
